\section{Background and Related Work}
\label{sec:background}

\subsection{Federated IoT Anomaly Detection and Training-Side Defenses}

Federated intrusion and anomaly detection for IoT has moved beyond plain FedAvg training~\cite{Belenguer2025FLIDSReview,Khang2026P4P,Quyen2025FedKDIDS,rey2022federated,Thein2024PFLIDS}: these studies improve training-time collaboration and robustness metrics, while thresholded autoencoder pipelines still rely on a benign post-training calibration step.
The N-BaIoT benchmark~\cite{meidan2018nbaiot} provides labeled benign and Mirai/BASHLITE traffic for nine physical devices and is also used in FL-IDS evaluation~\cite{Quyen2025FedKDIDS}.

The cited defenses act on the training phase.
Byzantine-robust rules, including Krum~\cite{blanchard2017machine}, Trimmed-Mean, and Median~\cite{yin2018byzantine}, filter or downweight outlier gradients before aggregation, and FLTrust~\cite{cao2020fltrust} anchors aggregation to a trusted server dataset.
Recent FL-IDS work extends this training-side view with personalized poisoning defenses, knowledge-distillation verifiers, and probe-guided anti-poisoning pipelines under non-IID IoT data~\cite{Khang2026P4P,Quyen2025FedKDIDS,Thein2024PFLIDS}.
These defenses inspect training data, gradients, model updates, or aggregation inputs; none is specified to observe the later client-local calibration buffer or designed for the post-training calibration channel considered here---a statement about the cited defenses, not a claim that no calibration-stage defense could be built.

\subsection{Threshold Calibration and Adjacent Work}

Kloft \& Laskov~\cite{kloft2010online} study online poisoning of centralized anomaly detection through a continuous training stream, rather than one-time post-training calibration.
FL model poisoning~\cite{bagdasaryan2020backdoor} likewise targets the gradient and weight surface.
Within federated anomaly detection, Laridi et al.~\cite{laridi2024} propose summary-statistics-based thresholding without adversarial contamination, and S\'aez-de-C\'amara et al.~\cite{saezdecamara2023} study clustered IoT detection. Calibration contamination has been studied elsewhere: Li et al.~\cite{li2024conformalpoison} attack conformal prediction, while Clarkson et al.~\cite{clarkson2024contamination} analyze split conformal prediction with contaminated calibration scores and a robust adjustment. These studies address different settings but establish that adversarial calibration has already been investigated outside federated anomaly detection.
The \GLOBAL, \LOCAL, and \CLUSTER policies evaluated here, defined in Section~\ref{sec:method}, follow the partitioning structures of that literature~\cite{ghosh2020efficient,laridi2024,sattler2020clustered}.

\begin{sloppypar}
The contribution here is a controlled comparison of policy-dependent victim and non-victim effects in one federated anomaly-detection setting. It does not establish a new poisoning taxonomy or traffic-level attack.
Section~\ref{sec:threat} states the adversary model.
\end{sloppypar}
